\documentclass[11pt]{article}
\usepackage[margin=1in]{geometry}
\usepackage{amsmath,amssymb,amsthm}
\usepackage{graphicx}
\usepackage[hidelinks]{hyperref}

\newtheorem{theorem}{Theorem}
\emergencystretch=2em
\newtheorem{remark}{Remark}
\DeclareMathOperator{\spanop}{span}

\begin{document}
\pagestyle{empty}

\begin{center}
{\Large Transverse local stability of the Gaussian propeller optimizer}\\[0.4em]
{\small Run \texttt{fa\_banach\_001}, arXiv:0807.4626, partial result}
\end{center}

\section*{Source problem}

Khot and Naor define an extremal Gaussian partition constant \(C(k)\) and
prove the equivalent Gaussian-width formula
\[
C(k)=\sup
\frac{\bigl(\mathbb E\max_{1\le j\le k}g_j\bigr)^2}
{\sum_{j=1}^k\mathbb E g_j^2},
\]
where the supremum is over centered Gaussian \(k\)-vectors. They ask whether
\[
C(k)\le C(3)=\frac{9}{8\pi}\qquad(k\ge2).
\]
Equivalently, does the global optimizer always use only the three blades of a
planar \(120^\circ\) propeller? The question and Gaussian formulation occur on
page 20 of \cite{KN}.

\begin{center}
\includegraphics[width=0.96\linewidth]{figures/open_problem_crop.png}
\end{center}

\section*{Partial result}

We prove strict local maximality against a natural genuinely
higher-dimensional class of perturbations: adjoining any fixed finite cloud
of small coefficient vectors perpendicular to the propeller plane.  The
vectors are centered across coordinates, which is the optimal translation
normalization for the Gaussian quotient.

\section*{Proof intuition}

Write \(M\) for the maximum of the three propeller coordinates. It dominates
a fixed multiple of the planar Gaussian radius, so
\(\mathbb P(M\le t)=O(t^2)\).  A new transverse coordinate of size
\(\varepsilon\) can beat the propeller only on this \(O(\varepsilon^2)\)
small-ball event, and its payoff there is \(O(\varepsilon)\). Thus all the new
coordinates improve the expected maximum by only \(O(\varepsilon^3)\).
Centering the enlarged coefficient family, however, raises its total variance
by a strictly positive \(O(\varepsilon^2)\) term. The quadratic loss therefore
beats the cubic gain for small nonzero \(\varepsilon\).

This does not prove the global propeller conjecture: perturbations with
longitudinal components, or configurations far from the propeller, remain
uncontrolled.

\begin{theorem}[Transverse local stability]\label{thm:main}
Let \(E\perp F\) be finite-dimensional real Hilbert spaces. Let
\(u_1,u_2,u_3\in E\) satisfy
\[
\|u_i\|=r>0,
\qquad
\langle u_i,u_j\rangle=-\frac{r^2}{2}\quad(i\ne j),
\]
and let \(w_1,\ldots,w_m\in F\), \(m\ge1\), be fixed and not all zero.
Put \(k=m+3\) and
\[
\bar w=\frac1k\sum_{j=1}^m w_j.
\]
For \(\varepsilon\ge0\), form the centered coefficient family
\[
v_i(\varepsilon)=u_i-\varepsilon\bar w\quad(1\le i\le3),
\qquad
v_{3+j}(\varepsilon)=\varepsilon(w_j-\bar w)
\quad(1\le j\le m).
\]
If \(G\) is standard Gaussian on \(E\oplus F\), define
\[
Q(\varepsilon)=
\frac{\bigl(\mathbb E\max_{1\le i\le k}
\langle G,v_i(\varepsilon)\rangle\bigr)^2}
{\sum_{i=1}^k\|v_i(\varepsilon)\|^2}.
\]
Then \(Q(0)=9/(8\pi)\), and there is
\(\varepsilon_0=\varepsilon_0(u_1,u_2,u_3,w_1,\ldots,w_m)>0\) such that
\[
Q(\varepsilon)<Q(0)=\frac{9}{8\pi}
\qquad(0<\varepsilon<\varepsilon_0).
\]
\end{theorem}

\begin{proof}
The hypotheses on the \(u_i\) imply \(u_1+u_2+u_3=0\), and the definition of
\(\bar w\) gives \(\sum_{i=1}^k v_i(\varepsilon)=0\).
Write \(G=X+Z\), where \(X\in E\) and \(Z\in F\) are independent standard
Gaussians, and set
\[
M=\max_{1\le i\le3}\langle X,u_i\rangle,
\qquad
N=\max_{1\le j\le m}\langle Z,w_j\rangle.
\]
Subtracting the same vector \(\varepsilon\bar w\) from every coefficient adds
the common centered Gaussian
\(-\varepsilon\langle Z,\bar w\rangle\) to their maximum.  Its expectation is
zero, hence
\begin{equation}\label{eq:numerator}
\mathbb E\max_{1\le i\le k}\langle G,v_i(\varepsilon)\rangle
=\mathbb E\max\{M,\varepsilon N\}.
\end{equation}

Put \(S=\sum_{i=1}^3\|u_i\|^2=3r^2\). Orthogonality of \(E\) and \(F\)
gives the exact denominator identity
\begin{equation}\label{eq:denominator}
\sum_{i=1}^k\|v_i(\varepsilon)\|^2=S+A\varepsilon^2,
\qquad
A=\sum_{j=1}^m\|w_j\|^2-\frac1k
\left\|\sum_{j=1}^m w_j\right\|^2.
\end{equation}
Since \(k=m+3\), Cauchy--Schwarz yields
\[
A\ge \left(1-\frac{m}{k}\right)\sum_{j=1}^m\|w_j\|^2
=\frac3k\sum_{j=1}^m\|w_j\|^2>0.
\]

Let \(R=\|X\|\). Among three directions separated by \(120^\circ\), one is
within \(60^\circ\) of every nonzero planar vector. Consequently
\begin{equation}\label{eq:radial}
M\ge \frac r2 R\ge0.
\end{equation}
As \(R\) is the radius of a standard two-dimensional Gaussian, for \(t\ge0\)
\begin{equation}\label{eq:smallball}
\mathbb P(M\le t)
\le \mathbb P(R\le 2t/r)
=1-e^{-2t^2/r^2}
\le \frac{2t^2}{r^2}.
\end{equation}

Let \(N_+=\max\{N,0\}\) and
\(\delta(\varepsilon)=\mathbb E[\max\{M,\varepsilon N\}-M]\).
Conditioning on \(Z\), and using independence together with
\eqref{eq:smallball}, gives
\begin{align*}
0\le\delta(\varepsilon)
&=\mathbb E(\varepsilon N-M)_+\\
&\le \mathbb E\!\left[
\varepsilon N_+\,\mathbb P(M\le\varepsilon N_+\mid Z)
\right]\\
&\le \frac{2\varepsilon^3}{r^2}\mathbb E N_+^3
=B\varepsilon^3,
\end{align*}
where \(B=2r^{-2}\mathbb E N_+^3\in(0,\infty)\).

Set \(W=\mathbb E M>0\). Equations \eqref{eq:numerator} and
\eqref{eq:denominator} imply
\[
Q(\varepsilon)\le
\frac{(W+B\varepsilon^3)^2}{S+A\varepsilon^2}.
\]
For example, choose
\[
\varepsilon_0=\min\left\{
\frac{WA}{8SB},
\left(\frac{W^2A}{4SB^2}\right)^{1/4}
\right\}>0.
\]
If \(0<\varepsilon<\varepsilon_0\), then
\[
2SWB\varepsilon+SB^2\varepsilon^4<W^2A,
\]
which is exactly
\(S(W+B\varepsilon^3)^2<W^2(S+A\varepsilon^2)\).  Hence
\(Q(\varepsilon)<W^2/S=Q(0)\).

For completeness, the last value can be computed directly.  In polar
coordinates, \(\mathbb ER=\sqrt{\pi/2}\), while the angular average of the
maximum of three cosines separated by \(120^\circ\) is
\(3\sqrt3/(2\pi)\).  Therefore
\[
W=r\sqrt{\frac\pi2}\frac{3\sqrt3}{2\pi}
=\frac{3\sqrt3}{2\sqrt{2\pi}}r,
\qquad
\frac{W^2}{S}=\frac{9}{8\pi}.
\]
\end{proof}

\section*{Verification notes}

The proof is analytic and has no numerical or computer-assisted dependency.
The two delicate points for human review are the centering identity
\eqref{eq:denominator} and the cubic gain estimate following
\eqref{eq:smallball}.  Both are displayed explicitly.

A bounded novelty search on 2026-07-22 covered the exact phrases
``propeller conjecture'', ``Gaussian moments'', ``local stability'', and
``orthogonal perturbation'', together with arXiv:0807.4626 and its principal
later low-dimensional result arXiv:1112.2993.  Heilman, Jagannath, and Naor
prove the full partition inequality in \(\mathbb R^3\), while explicitly
leaving higher dimensions open \cite{HJN}.  Searches also inspected the
related arXiv:1901.03934 and arXiv:2209.11216 abstracts.  No exact statement of
Theorem~\ref{thm:main} was found.  Novelty confidence is therefore moderate,
not definitive.

\enlargethispage{4\baselineskip}
\section*{Limitations and review recommendation}

This is a partial result, not a proof of \(C(k)\le C(3)\). It treats a fixed
finite transverse cloud scaled uniformly to zero.  It does not treat
longitudinal perturbations of the three active vertices, mixed
longitudinal--transverse new vertices, or configurations far from the
propeller.  Status: \textbf{candidate partial result, likely valid; human
review recommended}, with priority on the \(O(\varepsilon^3)\) gain versus
\(O(\varepsilon^2)\) variance mechanism.

\small
\begin{thebibliography}{9}
\bibitem{KN}
S. Khot and A. Naor,
\emph{Approximate kernel clustering},
arXiv:0807.4626, 2008.

\bibitem{HJN}
S. Heilman, A. Jagannath, and A. Naor,
\emph{Solution of the propeller conjecture in \(\mathbb R^3\)},
arXiv:1112.2993, 2011.
\end{thebibliography}
\normalsize

\end{document}
